\documentclass[12pt]{article}

\usepackage[a4paper, margin=1in]{geometry}
\usepackage{amsmath,amssymb,amsthm}
\usepackage{mathtools}
\usepackage{mathrsfs}
\usepackage{enumitem}
\usepackage{xcolor}
\usepackage[hidelinks]{hyperref}
\usepackage{fontspec}

\defaultfontfeatures{Ligatures=TeX}
\IfFontExistsTF{Times New Roman}
  {\setmainfont{Times New Roman}}
  {\setmainfont{TeX Gyre Termes}}
\IfFontExistsTF{Menlo}
  {\setmonofont{Menlo}}
  {\setmonofont{Latin Modern Mono}}

\newtheorem{theorem}{Theorem}[section]
\newtheorem{proposition}[theorem]{Proposition}
\newtheorem{lemma}[theorem]{Lemma}
\newtheorem{corollary}[theorem]{Corollary}
\theoremstyle{definition}
\newtheorem{definition}[theorem]{Definition}
\theoremstyle{remark}
\newtheorem{remark}[theorem]{Remark}

\newcommand{\tightlist}{%
  \setlength{\itemsep}{0pt}\setlength{\parskip}{0pt}}

\title{%
  \textbf{Foundation III: The Perfuming Loop}\\
  \textbf{A Conditional One-Step Source-Level Feedback Program}%
}

\author{
  Sidong Liu, PhD \\
  \small iBioStratix Ltd \\
  \small \texttt{sidongliu@hotmail.com}
}

\date{April 2026}

\begin{document}
\emergencystretch=2em
\raggedbottom
\hbadness=5000
\vbadness=5000

\maketitle

\begin{abstract}
Foundation~I~\cite{FoundationI} constructs a conditional single-source
chain from a source datum \(S_{\mathrm{source}}\) to an admissible
manifest pair \((\mathcal A_c,\omega)\). Foundation~II~\cite{FoundationII}
extends this to two sources, isolating a shared manifest sector. Both
chains are one-directional: source data produce manifest sectors, but
the manifest sectors do not feed back to the source. The present paper,
Foundation~III, addresses the return arrow.

T-DOME~III~\cite{TDOMEIII} establishes a calibration loop at the ego
level: the observer's prediction-error residuals drive a natural-gradient
update of the reference-frame parameter~\(\sigma\), with tracking
stability guaranteed by a Lyapunov bound. Paper~III of the
Consciousness Trilogy~\cite{PaperIII} proves that the resulting
self-description deficit \(\varepsilon_{\mathrm{loc}}(n)>0\) is
non-vanishing. Neither result closes the loop at the source level: the
ego adjusts its frame, but the decoherence semigroup
\(\{\mathcal E_t\}\) that generated the manifest sector in the first
place is left untouched.

Foundation~III defines a \emph{source-level feedback map} that lifts the
ego-level calibration signal to a controlled perturbation of the
decoherence semigroup. Under named regularity and coupling conditions,
we prove a conditional one-step closed-loop foothold with three
components. First, the feedback-modified semigroup remains a valid
source datum in the pointwise \(\sigma\)-weak Heisenberg-picture sense
used by Foundation~I v2. Second, the
self-description deficit \(\varepsilon_{\mathrm{loc}}\) under the
closed loop satisfies a modified bound that depends on the feedback
coupling strength. Third, the one-directional Foundation~I chain
upgrades only to a conditional one-step closed loop:
\[
\begin{aligned}
S_{\mathrm{source}}
&\longrightarrow
\mathcal O_B
\longrightarrow
\text{calibration signal}
\longrightarrow
S_{\mathrm{source}}[\text{updated}]\\
&\longrightarrow
\mathcal O_B[\text{updated}]
\longrightarrow
(\mathcal A_c,\omega)[\text{updated}].
\end{aligned}
\]
The paper does not claim that the loop eliminates the self-description
deficit. It proves only that the loop is well-posed and that the deficit
transforms in a controlled way. The question of whether iterated
perfuming converges to a fixed point, and what such a fixed point would
mean structurally, is left as an open problem.
\end{abstract}

\clearpage
\tableofcontents
\clearpage

\section{Introduction}\label{sec:intro}

Foundation~I~\cite{FoundationI} solved the first source-side bridge in
the Consciousness Series. Starting from a source datum
\[
S_{\mathrm{source}}
=
(\mathcal H_U,\mathcal A_U,\Omega,\{\mathcal E_t\}_{t\ge 0}),
\]
it isolated a conditional chain
\[
S_{\mathrm{source}}
\to
\mathcal O_A
\to
\mathcal O_B
\to
(\mathcal A_c,\omega),
\]
where \(\mathcal O_A\) is a persistent projected record process,
\(\mathcal O_B\) is a budget-forced foreground frame, and
\((\mathcal A_c,\omega)\) is the manifest accessible pair. The chain is
one-directional: source data produce a manifest sector, and the series
then hands that sector to Paper~I of the Consciousness Trilogy
\cite{PaperI} for the geometry--modular-flow theorem and to
Paper~II~\cite{PaperII} for the higher-dimensional consistency analysis.

Foundation~II~\cite{FoundationII} extended that source-side analysis to
the first multi-source setting. There the main result was again
one-directional. Given two source chains sharing one ambient
Type~III\(_1\) substrate, one obtains a conditional inclusion theorem
for a shared admissible manifest sector. What remained absent in both
papers was any formal return arrow from manifest calibration back to the
source dynamics that generated the manifest sector in the first place.
Foundation~II enters the present paper only as programme boundary
calibration for the later multi-source feedback problem; no theorem
below imports a Foundation~II result.

The missing input now comes from two already published directions.
T-DOME~III~\cite{TDOMEIII} supplied a one-step calibration signal at the
ego level: prediction-residual statistics determine a natural-gradient
update \(\delta\sigma\), controlled by a Lyapunov tracking estimate and
penalized by the ego-rigidity term
\(\mathcal I_{\mathrm{ego}}\). Paper~III of the Consciousness Trilogy
\cite{PaperIII} supplied a quantitative obstruction on the manifest
side: the local self-description deficit
\(\varepsilon_{\mathrm{loc}}(n)\) is bounded away from zero for
locally bounded \(n\)-dimensional self-models. Neither result, by
itself, modifies the source semigroup \(\{\mathcal E_t\}\).

The present paper addresses exactly that gap. Its aim is not to prove a
global convergence theorem for iterative self-modification or a general
closed-loop theorem. It does not claim that the self-description deficit
is eliminated by feedback, and it does not claim that one can derive a
feedback lift from bare multi-scale dynamics without additional
structure. Its narrower goal is to identify an honest one-step
theorem-level foothold for \emph{source-level feedback from manifest
calibration}.

The main design choice is to avoid a formal expression of the type
\[
\mathcal E_t \mapsto \mathcal E_t + \eta\,\Gamma(\delta\sigma,\mathcal E_t)
\]
as the \emph{definition} of the modified dynamics. Such a first-order
expression is useful heuristically, but it does not by itself preserve
the semigroup law or complete positivity. Instead, the paper defines a
feedback lift through an exact auxiliary semigroup
\(\{\mathscr F_r^{[\delta\sigma]}\}_{r\ge 0}\) of normal completely
positive unital maps and then sets
\[
\mathcal E_t^{[\eta]}
:=
\mathscr F_{\eta t}^{[\delta\sigma]}\circ\mathcal E_t.
\]
Once a commutation condition is imposed, the semigroup law is automatic;
the first-order feedback map \(\Gamma\) then reappears as the
\(\eta\downarrow 0\) tangent of this exact construction.

This paper therefore contributes four one-step theorem-level items.
\begin{enumerate}
\item We define a \emph{source-level feedback lift} from an admissible
  calibration signal \(\delta\sigma\) to an exact Heisenberg-picture
  semigroup perturbation on the ambient Type~III\(_1\) algebra.
\item We prove a preservation theorem: under named regularity and
  commutation inputs, the modified family
  \(\{\mathcal E_t^{[\eta]}\}_{t\ge0}\) is again a pointwise
  \(\sigma\)-weakly continuous semigroup of normal completely positive
  unital maps on \(\mathcal A_U\), so the ambient datum remains a valid
  source datum.
\item We prove a \emph{modified deficit-floor theorem}: if the chosen
  Foundation~I chain is feedback-stable and the local Paper~III Fisher
  constants remain controlled, then the Paper~III lower floor deforms by
  an explicit small-coupling factor \(f(\eta)\) with \(f(0)=1\).
\item We package the result as a \emph{one-step closed-loop datum},
  showing that the chain
  \[
  S_{\mathrm{source}}
  \to
  \mathcal O_B
  \to
  \delta\sigma
  \to
  S_{\mathrm{source}}^{[\eta]}
  \to
  \mathcal O_B^{[\eta]}
  \to
  (\mathcal A_c^{[\eta]},\omega^{[\eta]})
  \]
  is mathematically well-posed under named inputs.
\end{enumerate}

The scope is deliberately narrow. The paper does not solve the
constructive origin of the calibration signal, which is imported from
T-DOME~III. It does not solve the full persistence problem for the
Foundation~I chain under arbitrary perturbations, which remains a
conditional interface. It does not derive explicit optimal coupling
windows or fixed points for iterated perfuming. Those are precisely the
open problems recorded later.

The paper also keeps the same Type~III discipline as Foundation~II. No
tensor-product factorization is assumed. All dynamics are written in the
Heisenberg picture on a common ambient algebra, and continuity means
pointwise \(\sigma\)-weak continuity rather than norm continuity. This
is not cosmetic: the whole point of the paper is to modify the source
semigroup without importing an external finite-dimensional controller.

The paper is organized as follows. Section~\ref{sec:feedback} defines
the feedback lift, the regularity package, and the modified source
datum. Section~\ref{sec:preservation} proves the source-datum
preservation theorem. Section~\ref{sec:deficit} states the extra
comparison inputs needed to pass through Foundation~I and then proves
the modified deficit-floor theorem. Section~\ref{sec:closed} packages
the resulting one-step closed loop. Section~\ref{sec:open} records the
open problems. Section~\ref{sec:ledger} collects the named conditional
inputs. Appendix~\ref{app:interfaces} restates the imported interfaces
from Foundation~I, T-DOME~III, and Paper~III in theorem-ready form.

\section{Feedback Lift and Regularity Conditions}\label{sec:feedback}

The first task is to isolate the additional object that Foundation~III
needs and the exact conditions under which it is allowed to act on the
source semigroup.

\begin{definition}[Admissible calibration signal]\label{def:signal}
Fix a source datum \(S_{\mathrm{source}}\) together with a chosen
Foundation~I chain up to Stage~B, so that one has
\[
\mathcal O_B
=
(\mathcal P,\omega_{\mathcal P,t},\mathcal K_*(t,s),\mathfrak F,\pi_{\mathrm{ego}}).
\]
An \emph{admissible calibration signal} is a vector
\[
\delta\sigma \in T_\sigma\Sigma
\]
in the local calibration parameter space of T-DOME~III such that:
\begin{enumerate}
\item \(\delta\sigma\) is obtained from the natural-gradient calibration
  step imported from T-DOME~III;
\item \(\sigma+\delta\sigma\) remains inside the Lyapunov neighborhood on
  which the T-DOME~III tracking estimate is valid;
\item there exists a bound
  \(\|\delta\sigma\| \leq R_{\mathrm{cal}}\) for some finite
  calibration radius \(R_{\mathrm{cal}}>0\).
\end{enumerate}
\end{definition}

\begin{remark}[One-step rather than iterative calibration]
Definition~\ref{def:signal} is intentionally local. The paper treats a
single feedback step driven by one calibration signal
\(\delta\sigma\). An iterated sequence
\(\delta\sigma_0,\delta\sigma_1,\dots\) would require a separate fixed
point or long-time theory, and nothing in the present paper should be
read as supplying that theory.
\end{remark}

\begin{definition}[Feedback regularity package]\label{def:FB}
Let \(S_{\mathrm{source}}\) be a source datum and let
\(\delta\sigma\) be an admissible calibration signal. A
\emph{feedback regularity package} consists of a family
\[
\mathscr F^{[\delta\sigma]}
=
\{\mathscr F_r^{[\delta\sigma]}\}_{r\ge0}
\]
of normal maps on \(\mathcal A_U\) satisfying the following conditions.
\begin{enumerate}
\item \textbf{FB-A: CP source lift.}
  \(\{\mathscr F_r^{[\delta\sigma]}\}_{r\ge0}\) is a pointwise
  \(\sigma\)-weakly continuous semigroup of normal completely positive
  unital maps on \(\mathcal A_U\), with
  \[
  \mathscr F_0^{[\delta\sigma]}=\mathrm{id}_{\mathcal A_U},
  \qquad
  \mathscr F_r^{[0]}=\mathrm{id}_{\mathcal A_U}
  \quad
  \forall r\ge0.
  \]
\item \textbf{FB-B: small-signal differentiability.}
  For every \(A\in\mathcal A_U\), the map
  \(r\mapsto \mathscr F_r^{[\delta\sigma]}(A)\) is
  \(\sigma\)-weakly differentiable at \(r=0\). We denote the
  derivative by
  \[
  \mathfrak G_{\delta\sigma}(A)
  :=
  \left.
  \frac{d}{dr}\mathscr F_r^{[\delta\sigma]}(A)
  \right|_{r=0},
  \]
  and record the normalization
  \(\mathfrak G_{\delta\sigma}(\mathbf 1)=0\), which follows from the
  unitality in FB-A because
  \(\mathscr F_r^{[\delta\sigma]}(\mathbf 1)=\mathbf 1\) for all
  \(r\ge0\).
\item \textbf{FB-C: chain compatibility.}
  The feedback lift commutes with the ambient source semigroup:
  \[
  \mathscr F_r^{[\delta\sigma]}\mathcal E_t
  =
  \mathcal E_t\mathscr F_r^{[\delta\sigma]}
  \qquad
  \forall r,t\ge0.
  \]
  In addition, for the chosen Foundation~I chain, the maps
  \(\mathscr F_r^{[\delta\sigma]}\) preserve the named structural inputs
  that are required later to reconstruct Stage~A, Stage~B, and the
  manifest pair.
\end{enumerate}
\end{definition}

\begin{remark}[Why the feedback acts through an exact semigroup]
The exact feedback semigroup in Definition~\ref{def:FB} is not an
ornament. A first-order perturbation
\(\mathcal E_t+\eta\,\Gamma(\delta\sigma,\mathcal E_t)\) does not
automatically preserve the semigroup property or complete positivity.
The exact semigroup lift is the cleanest way to keep those properties in
the Heisenberg-picture Type~III setting. The first-order map
\(\Gamma\) reappears only after the exact object has been defined.
\end{remark}

\begin{remark}[Meaning of the strict commutation input]
The commutation requirement in FB-C is strong and should be read
carefully. It does \emph{not} say that feedback is dynamically trivial.
For \(\eta>0\),
\[
\mathcal E_t^{[\eta]}
=
\mathscr F_{\eta t}^{[\delta\sigma]}\circ \mathcal E_t
\]
generally differs from \(\mathcal E_t\) unless the calibration signal is
zero or the coupling is turned off. The role of strict commutation is
only to keep the algebra-level time ordering unchanged, so that the
modified dynamics remains an exact semigroup without entering a more
delicate generator-level perturbation problem. Relaxing FB-C to an
approximate-commutation condition is left open.
\end{remark}

\begin{definition}[Small-coupling regime and modified semigroup]
\label{def:modified}
Fix a feedback regularity package
\(\mathscr F^{[\delta\sigma]}\). A coupling strength \(\eta\ge0\) is in
the \emph{small-coupling regime} if \(0\le \eta < \eta_*\) for some
window \(\eta_*>0\) on which the later comparison inputs remain valid.
For such \(\eta\), define the \emph{feedback-modified semigroup} by
\[
\mathcal E_t^{[\eta]}
:=
\mathscr F_{\eta t}^{[\delta\sigma]}\circ \mathcal E_t
\qquad
(t\ge0).
\]
The corresponding \emph{modified source datum} is
\[
S_{\mathrm{source}}^{[\eta]}
:=
(\mathcal H_U,\mathcal A_U,\Omega,\{\mathcal E_t^{[\eta]}\}_{t\ge0}).
\]
\end{definition}

\begin{remark}[Heisenberg-picture convention]
As in Foundation~II, all semigroups in Definition~\ref{def:modified}
act on observables. Pointwise continuity therefore means
\(\sigma\)-weak continuity on \(\mathcal A_U\), not norm continuity.
Nothing in the present paper uses or claims a norm-continuous
generator-level perturbation theorem.
\end{remark}

\begin{definition}[First-order feedback map]\label{def:Gamma}
Under FB-B, define the \emph{first-order feedback map} by
\[
\Gamma(\delta\sigma,\mathcal E_t)(A)
:=
t\,\mathfrak G_{\delta\sigma}(\mathcal E_t(A)),
\qquad
A\in\mathcal A_U.
\]
Equivalently, for each fixed \(t\ge0\),
\[
\mathcal E_t^{[\eta]}(A)
=
\mathcal E_t(A)
+
\eta\,\Gamma(\delta\sigma,\mathcal E_t)(A)
+
o(\eta)
\]
in the \(\sigma\)-weak topology as \(\eta\downarrow 0\).
\end{definition}

\begin{proof}[Verification of the expansion]
By Definition~\ref{def:modified},
\[
\mathcal E_t^{[\eta]}(A)
=
\mathscr F_{\eta t}^{[\delta\sigma]}(\mathcal E_t(A)).
\]
Applying the \(r=0\) derivative from FB-B at \(r=\eta t\) gives
\[
\mathscr F_{\eta t}^{[\delta\sigma]}(\mathcal E_t(A))
=
\mathcal E_t(A)
+
\eta t\,\mathfrak G_{\delta\sigma}(\mathcal E_t(A))
+
o(\eta),
\]
which is exactly the stated formula.
\end{proof}

\begin{remark}[What Definition~\ref{def:Gamma} does not prove]
The first-order map \(\Gamma\) is merely the tangent of the exact
feedback lift. Definition~\ref{def:Gamma} does not claim that an
arbitrary bounded normal map \(\Gamma\) integrates to a completely
positive semigroup, nor does it characterize the most general
generator-level perturbations compatible with Type~III dynamics. Those
questions are left open.
\end{remark}

\section{Source-Datum Preservation}\label{sec:preservation}

We now prove the first of the paper's three main results.

\begin{lemma}[Semigroup law]\label{lem:semigroup}
Assume FB-A--FB-C. Then the family
\(\{\mathcal E_t^{[\eta]}\}_{t\ge0}\) from
Definition~\ref{def:modified} satisfies
\[
\mathcal E_t^{[\eta]}\mathcal E_s^{[\eta]}
=
\mathcal E_{t+s}^{[\eta]}
\qquad
\forall s,t\ge0.
\]
\end{lemma}

\begin{proof}
For \(A\in\mathcal A_U\),
\begin{align*}
\mathcal E_t^{[\eta]}\mathcal E_s^{[\eta]}(A)
&=
\mathscr F_{\eta t}^{[\delta\sigma]}\,
\mathcal E_t
\Bigl(
\mathscr F_{\eta s}^{[\delta\sigma]}\,
\mathcal E_s(A)
\Bigr) \\
&=
\mathscr F_{\eta t}^{[\delta\sigma]}
\mathscr F_{\eta s}^{[\delta\sigma]}
\mathcal E_t\mathcal E_s(A),
\end{align*}
where FB-C is used to commute
\(\mathcal E_t\) past \(\mathscr F_{\eta s}^{[\delta\sigma]}\). Since
\(\{\mathscr F_r^{[\delta\sigma]}\}\) and \(\{\mathcal E_t\}\) are both
semigroups,
\[
\mathscr F_{\eta t}^{[\delta\sigma]}
\mathscr F_{\eta s}^{[\delta\sigma]}
\mathcal E_t\mathcal E_s
=
\mathscr F_{\eta(t+s)}^{[\delta\sigma]}
\mathcal E_{t+s}
=
\mathcal E_{t+s}^{[\eta]}.
\]
\end{proof}

\begin{lemma}[Pointwise \(\sigma\)-weak continuity and positivity]
\label{lem:continuity}
Assume FB-A--FB-C. Then for every \(A\in\mathcal A_U\),
\[
t\longmapsto \mathcal E_t^{[\eta]}(A)
\]
is pointwise \(\sigma\)-weakly continuous. Moreover, each
\(\mathcal E_t^{[\eta]}\) is normal, completely positive, and unital.
\end{lemma}

\begin{proof}
Fix \(A\in\mathcal A_U\) and a normal functional
\(\varphi\in(\mathcal A_U)_*\). By Definition~\ref{def:modified},
\[
\varphi(\mathcal E_t^{[\eta]}(A))
=
\varphi\!\left(
\mathscr F_{\eta t}^{[\delta\sigma]}(\mathcal E_t(A))
\right).
\]
The map \(t\mapsto \mathcal E_t(A)\) is \(\sigma\)-weakly continuous by
the source-datum assumption, and
\(r\mapsto \mathscr F_r^{[\delta\sigma]}(B)\) is
\(\sigma\)-weakly continuous for every fixed \(B\in\mathcal A_U\) by
FB-A. To prove continuity of the composition, fix \(t_0\ge0\) and write
\begin{align*}
&\left|
\varphi\bigl(\mathscr F_{\eta t}^{[\delta\sigma]}(\mathcal E_t(A))\bigr)
-
\varphi\bigl(\mathscr F_{\eta t_0}^{[\delta\sigma]}(\mathcal E_{t_0}(A))\bigr)
\right| \\
&\qquad \le
\left|
\varphi\bigl(\mathscr F_{\eta t}^{[\delta\sigma]}(\mathcal E_t(A))\bigr)
-
\varphi\bigl(\mathscr F_{\eta t}^{[\delta\sigma]}(\mathcal E_{t_0}(A))\bigr)
\right| \\
&\qquad\quad +
\left|
\varphi\bigl(\mathscr F_{\eta t}^{[\delta\sigma]}(\mathcal E_{t_0}(A))\bigr)
-
\varphi\bigl(\mathscr F_{\eta t_0}^{[\delta\sigma]}(\mathcal E_{t_0}(A))\bigr)
\right|.
\end{align*}
For the second term, \(\mathcal E_{t_0}(A)\) is fixed, so pointwise
\(\sigma\)-weak continuity of
\(\{\mathscr F_r^{[\delta\sigma]}\}_{r\ge0}\) gives convergence to
zero as \(t\to t_0\).

For the first term, use the predual semigroup
\(\{\mathscr F_{r,*}^{[\delta\sigma]}\}_{r\ge0}\). Since
\(\{\mathscr F_r^{[\delta\sigma]}\}\) is a pointwise
\(\sigma\)-weakly continuous semigroup of normal maps, the predual
family is norm-continuous on \((\mathcal A_U)_*\). Thus
\begin{align*}
&\left|
\varphi\bigl(\mathscr F_{\eta t}^{[\delta\sigma]}(\mathcal E_t(A))\bigr)
-
\varphi\bigl(\mathscr F_{\eta t}^{[\delta\sigma]}(\mathcal E_{t_0}(A))\bigr)
\right| \\
&\qquad =
\left|
\mathscr F_{\eta t,*}^{[\delta\sigma]}(\varphi)
\bigl(\mathcal E_t(A)-\mathcal E_{t_0}(A)\bigr)
\right| \\
&\qquad \le
\left|
\bigl(\mathscr F_{\eta t,*}^{[\delta\sigma]}(\varphi)
-
\mathscr F_{\eta t_0,*}^{[\delta\sigma]}(\varphi)\bigr)
\bigl(\mathcal E_t(A)-\mathcal E_{t_0}(A)\bigr)
\right| \\
&\qquad\quad +
\left|
\mathscr F_{\eta t_0,*}^{[\delta\sigma]}(\varphi)
\bigl(\mathcal E_t(A)-\mathcal E_{t_0}(A)\bigr)
\right|.
\end{align*}
The first summand tends to zero because
\(\mathscr F_{\eta t,*}^{[\delta\sigma]}(\varphi)\to
\mathscr F_{\eta t_0,*}^{[\delta\sigma]}(\varphi)\) in predual norm,
while \(\mathcal E_t(A)-\mathcal E_{t_0}(A)\) is norm-bounded by
\(2\|A\|\). The second summand tends to zero because
\(\mathcal E_t(A)\to\mathcal E_{t_0}(A)\) \(\sigma\)-weakly and
\(\mathscr F_{\eta t_0,*}^{[\delta\sigma]}(\varphi)\) is a fixed normal
functional. Hence the first term also tends to zero.

Hence \(\{\mathcal E_t^{[\eta]}\}\) is pointwise \(\sigma\)-weakly
continuous.

Normality and complete positivity are preserved under composition of
normal completely positive maps, and unitality is preserved under
composition of unital maps. Since both
\(\mathscr F_{\eta t}^{[\delta\sigma]}\) and \(\mathcal E_t\) have
those properties, so does \(\mathcal E_t^{[\eta]}\).
\end{proof}

\begin{theorem}[Source-datum preservation]\label{thm:preserve}
Let
\[
S_{\mathrm{source}}
=
(\mathcal H_U,\mathcal A_U,\Omega,\{\mathcal E_t\}_{t\ge0})
\]
be a source datum in the sense imported from Foundation~I, let
\(\delta\sigma\) be an admissible calibration signal, and assume that a
feedback regularity package FB-A--FB-C is given. Then for every
\(\eta\in[0,\eta_*)\), the modified datum
\[
S_{\mathrm{source}}^{[\eta]}
=
(\mathcal H_U,\mathcal A_U,\Omega,\{\mathcal E_t^{[\eta]}\}_{t\ge0})
\]
is again a valid source datum: the family
\(\{\mathcal E_t^{[\eta]}\}_{t\ge0}\) is a pointwise
\(\sigma\)-weakly continuous semigroup of normal completely positive
unital maps on \(\mathcal A_U\), and \(\Omega\) remains cyclic and
separating for \(\mathcal A_U\).
\end{theorem}

\begin{proof}
Lemma~\ref{lem:semigroup} gives the semigroup law and
Lemma~\ref{lem:continuity} gives pointwise \(\sigma\)-weak continuity,
normality, complete positivity, and unitality.

The ambient standard-form triple
\((\mathcal H_U,\mathcal A_U,\Omega)\) is not changed by the feedback
step. In particular, \(\Omega\) remains cyclic and separating for the
same von Neumann algebra \(\mathcal A_U\). Therefore
\(S_{\mathrm{source}}^{[\eta]}\) satisfies the same source-datum format
as \(S_{\mathrm{source}}\).
\end{proof}

\begin{remark}[What Theorem~\ref{thm:preserve} does not prove]
The theorem proves only that the feedback-modified dynamics remains in
the theorem domain of the Foundation~I source-datum definition. It does
not prove that the modified datum automatically reproduces the same
persistent projection, the same foreground frame, or the same manifest
algebra. Those are downstream interface questions and require additional
inputs.
\end{remark}

\begin{corollary}[Exact CP lift with first-order tangent]
\label{cor:firstorder}
Under the hypotheses of Theorem~\ref{thm:preserve}, the map
\(\Gamma(\delta\sigma,\mathcal E_t)\) from
Definition~\ref{def:Gamma} is the \(\eta\downarrow 0\) tangent of an
exact Heisenberg-picture completely positive perturbation.
\end{corollary}

\begin{proof}
This is immediate from Theorem~\ref{thm:preserve} together with the
first-order expansion of Definition~\ref{def:Gamma}.
\end{proof}

\section{Modified Deficit Floor}\label{sec:deficit}

The preservation theorem keeps the modified dynamics inside the
source-datum domain, but it does not yet produce a manifest pair. For
that we need the chosen Foundation~I chain to survive the perturbation
and the local Paper~III constants to remain controlled.

\begin{definition}[Deficit-comparison package]\label{def:DFA}
Fix a source datum \(S_{\mathrm{source}}\), a chosen Foundation~I chain
from \(S_{\mathrm{source}}\) to \((\mathcal A_c,\omega)\), and a
feedback-modified datum \(S_{\mathrm{source}}^{[\eta]}\). We say that
\emph{DF-A} holds on the coupling window \([0,\eta_*)\) if:
\begin{enumerate}
\item \textbf{DF-A(1): feedback-stable Foundation~I chain.}
  For every \(\eta\in[0,\eta_*)\), the named conditional inputs used in
  the chosen Foundation~I chain remain available for
  \(S_{\mathrm{source}}^{[\eta]}\). Hence one obtains a modified manifest
  pair
  \[
  (\mathcal A_c^{[\eta]},\omega^{[\eta]}).
  \]
\item \textbf{DF-A(2): Fisher control.}
  There exist continuous functions
  \(\alpha_F,\beta_{\mathrm{ego}}:[0,\eta_*)\to(0,\infty)\) with
  \[
  \alpha_F(0)=\beta_{\mathrm{ego}}(0)=1
  \]
  such that
  \[
  \mathcal I_F^{[\eta]}
  \ge
  \alpha_F(\eta)\,\mathcal I_F,
  \qquad
  \mathcal I_{\mathrm{ego}}^{[\eta]}
  \le
  \beta_{\mathrm{ego}}(\eta)\,\mathcal I_{\mathrm{ego}},
  \]
  where \(\mathcal I_F\) and
  \(\mathcal I_{\mathrm{ego}}\) are the baseline manifest-side Fisher
  sensitivity and ego-rigidity constants, and
  \(\mathcal I_F^{[\eta]}\),
  \(\mathcal I_{\mathrm{ego}}^{[\eta]}\) are the corresponding modified
  quantities.
\item \textbf{DF-A(3): local Paper~III control.}
  The modified manifest pair
  \((\mathcal A_c^{[\eta]},\omega^{[\eta]})\) satisfies the local
  Paper~III hypotheses on a controlled neighborhood, with quadratic-loss
  parameter \(\rho(\eta)\in[0,\tfrac12)\) continuous in \(\eta\), and
  baseline value \(\rho(0)=:\rho_0\).
\end{enumerate}
\end{definition}

\begin{proposition}[Passage from the modified source datum to the modified manifest pair]
\label{prop:chain}
Assume the hypotheses of Theorem~\ref{thm:preserve} and DF-A(1). Then
for every \(\eta\in[0,\eta_*)\), the modified source datum
\(S_{\mathrm{source}}^{[\eta]}\) admits a conditional Foundation~I chain
\[
S_{\mathrm{source}}^{[\eta]}
\to
\mathcal O_A^{[\eta]}
\to
\mathcal O_B^{[\eta]}
\to
(\mathcal A_c^{[\eta]},\omega^{[\eta]}).
\]
\end{proposition}

\begin{proof}
Theorem~\ref{thm:preserve} puts
\(S_{\mathrm{source}}^{[\eta]}\) inside the theorem domain of the
Foundation~I source-datum definition. DF-A(1) states exactly that the
named conditional inputs used by the chosen Foundation~I chain remain
available after the feedback perturbation. The Foundation~I conditional
theorem can therefore be re-run with the modified datum, yielding the
stated outputs.
\end{proof}

\begin{definition}[Baseline and modified local floors]
\label{def:floor}
Under DF-A, define the baseline Paper~III local floor by
\[
\underline{\varepsilon}_{\mathrm{loc}}(n)
:=
\frac{(1-2\rho_0)\,\mathcal I_F}
{2\bigl(\mathcal I_F+\mathcal I_{\mathrm{ego}}\bigr)},
\]
and the modified local floor by
\[
\underline{\varepsilon}_{\mathrm{loc}}^{[\eta]}(n)
:=
\frac{(1-2\rho(\eta))\,\mathcal I_F^{[\eta]}}
{2\bigl(\mathcal I_F^{[\eta]}+\mathcal I_{\mathrm{ego}}^{[\eta]}\bigr)}.
\]
These are the theorem-level lower-floor quantities imported from
Paper~III for the baseline and modified manifest pairs.
\end{definition}

\begin{theorem}[Modified local deficit floor]\label{thm:deficit}
Assume the hypotheses of Proposition~\ref{prop:chain} and the full
package DF-A(1)--DF-A(3). Then for every \(\eta\in[0,\eta_*)\),
\[
\underline{\varepsilon}_{\mathrm{loc}}^{[\eta]}(n)
\ge
f(\eta)\,
\underline{\varepsilon}_{\mathrm{loc}}(n),
\]
where
\[
f(\eta)
:=
\frac{1-2\rho(\eta)}{1-2\rho_0}\,
\frac{\alpha_F(\eta)
\bigl(\mathcal I_F+\mathcal I_{\mathrm{ego}}\bigr)}
{\alpha_F(\eta)\mathcal I_F
+\beta_{\mathrm{ego}}(\eta)\mathcal I_{\mathrm{ego}}}.
\]
In particular, \(f(0)=1\), and by continuity \(f(\eta)>0\) for all
sufficiently small \(\eta\).
\end{theorem}

\begin{proof}
By Proposition~\ref{prop:chain}, the modified datum yields a manifest
pair \((\mathcal A_c^{[\eta]},\omega^{[\eta]})\). By DF-A(3), that pair
lies in the local theorem domain of Paper~III, so the Paper~III local
drift-floor theorem gives
\[
\varepsilon_{\mathrm{loc}}^{[\eta]}(n)
\ge
\underline{\varepsilon}_{\mathrm{loc}}^{[\eta]}(n)
=
\frac{(1-2\rho(\eta))\,\mathcal I_F^{[\eta]}}
{2\bigl(\mathcal I_F^{[\eta]}+\mathcal I_{\mathrm{ego}}^{[\eta]}\bigr)}.
\]
Using DF-A(2),
\[
\mathcal I_F^{[\eta]}
\ge
\alpha_F(\eta)\mathcal I_F,
\qquad
\mathcal I_{\mathrm{ego}}^{[\eta]}
\le
\beta_{\mathrm{ego}}(\eta)\mathcal I_{\mathrm{ego}},
\]
so
\[
\underline{\varepsilon}_{\mathrm{loc}}^{[\eta]}(n)
\ge
\frac{(1-2\rho(\eta))\,\alpha_F(\eta)\mathcal I_F}
{2\bigl(\alpha_F(\eta)\mathcal I_F
+\beta_{\mathrm{ego}}(\eta)\mathcal I_{\mathrm{ego}}\bigr)}.
\]
Now factor out the baseline floor from
Definition~\ref{def:floor}:
\begin{align*}
\underline{\varepsilon}_{\mathrm{loc}}^{[\eta]}(n)
&\ge
\frac{1-2\rho(\eta)}{1-2\rho_0}\,
\frac{\alpha_F(\eta)
\bigl(\mathcal I_F+\mathcal I_{\mathrm{ego}}\bigr)}
{\alpha_F(\eta)\mathcal I_F
+\beta_{\mathrm{ego}}(\eta)\mathcal I_{\mathrm{ego}}}
\cdot
\frac{(1-2\rho_0)\mathcal I_F}
{2(\mathcal I_F+\mathcal I_{\mathrm{ego}})} \\
&=
f(\eta)\,
\underline{\varepsilon}_{\mathrm{loc}}(n).
\end{align*}
The formula for \(f(0)\) follows from
\(\rho(0)=\rho_0\) and \(\alpha_F(0)=\beta_{\mathrm{ego}}(0)=1\).
Continuity then gives positivity for sufficiently small \(\eta\).
\end{proof}

\begin{remark}[No monotonicity claim for \(f(\eta)\)]
The theorem asserts only that \(f(0)=1\) and that \(f(\eta)\) remains
positive for sufficiently small coupling. No monotonicity statement is
claimed. Depending on the detailed response of the manifest Fisher
sensitivity and the ego-rigidity term, one may have \(f(\eta)<1\) or
\(f(\eta)>1\). In particular, a feedback step that increases
\(\mathcal I_F^{[\eta]}\) or decreases
\(\mathcal I_{\mathrm{ego}}^{[\eta]}\) can raise the lower floor rather
than lower it.
\end{remark}

\begin{corollary}[Comparison with the actual local deficit under saturation]
\label{cor:saturation}
Assume, in addition to the hypotheses of
Theorem~\ref{thm:deficit}, that the baseline Paper~III lower floor is
attained by an optimal locally bounded self-model, so that
\[
\varepsilon_{\mathrm{loc}}(n)
=
\underline{\varepsilon}_{\mathrm{loc}}(n).
\]
Then
\[
\varepsilon_{\mathrm{loc}}^{[\eta]}(n)
\ge
f(\eta)\,\varepsilon_{\mathrm{loc}}(n).
\]
\end{corollary}

\begin{proof}
By Theorem~\ref{thm:deficit},
\[
\varepsilon_{\mathrm{loc}}^{[\eta]}(n)
\ge
\underline{\varepsilon}_{\mathrm{loc}}^{[\eta]}(n)
\ge
f(\eta)\,\underline{\varepsilon}_{\mathrm{loc}}(n).
\]
Under the saturation hypothesis,
\(\underline{\varepsilon}_{\mathrm{loc}}(n)=
\varepsilon_{\mathrm{loc}}(n)\), which yields the claim.
\end{proof}

\begin{remark}[Status of the saturation hypothesis]
The saturation assumption in Corollary~\ref{cor:saturation} is a genuine
extra hypothesis. Paper~III proves the lower floor, but it does not in
general prove that the lower floor is attained by an optimal locally
bounded self-model. The corollary is therefore only an optional bridge
from the theorem-level lower-floor comparison to a comparison of the
actual minimax deficit.
\end{remark}

\begin{remark}[What Theorem~\ref{thm:deficit} does not prove]
The theorem gives a controlled deformation of the \emph{Paper~III lower
floor}. It does not prove that the true minimax deficit is monotone in
\(\eta\), it does not prove that feedback reduces the deficit, and it
does not prove that the deficit can be driven to zero. Those stronger
claims would require a full transport theory for optimal self-models,
not just control of the lower-floor constants.
\end{remark}

\begin{remark}[Why the comparison factor is the right theorem target]
The factor \(f(\eta)\) isolates exactly what feedback changes at theorem
level:
\begin{enumerate}
\item the local quadratic-loss constant \(\rho(\eta)\) from the
  Paper~III local entropy expansion;
\item the manifest Fisher sensitivity
  \(\mathcal I_F^{[\eta]}\);
\item the ego-rigidity contribution
  \(\mathcal I_{\mathrm{ego}}^{[\eta]}\).
\end{enumerate}
No stronger conclusion is claimed. In particular, the theorem does not
say that a better calibration signal necessarily increases
\(\mathcal I_F^{[\eta]}\) or decreases
\(\mathcal I_{\mathrm{ego}}^{[\eta]}\); it only says that if those
quantities remain controlled, then the local deficit floor remains
controlled as well.
\end{remark}

\section{Closed-Loop Well-Posedness}\label{sec:closed}

We can now state the paper's final theorem.

\begin{definition}[Closed-loop datum]\label{def:closed}
Under the hypotheses of Theorem~\ref{thm:deficit}, the
\emph{closed-loop datum at coupling \(\eta\)} is the tuple
\[
\mathfrak C_\eta
:=
\bigl(
S_{\mathrm{source}},
\mathcal O_B,
(\mathcal A_c,\omega),
\delta\sigma,
S_{\mathrm{source}}^{[\eta]},
\mathcal O_B^{[\eta]},
(\mathcal A_c^{[\eta]},\omega^{[\eta]})
\bigr),
\]
where \(\mathcal O_B\) and \((\mathcal A_c,\omega)\) are the baseline
Stage~B and Stage~C outputs of the chosen Foundation~I chain,
\(\delta\sigma\) is the admissible calibration signal extracted from the
retained Stage~B data, and
\(\mathcal O_B^{[\eta]}\) together with
\((\mathcal A_c^{[\eta]},\omega^{[\eta]})\) are the corresponding
modified outputs after the feedback step.
\end{definition}

\begin{theorem}[Closed-loop well-posedness]\label{thm:closed}
Let \(S_{\mathrm{source}}\) be a source datum, let
\(\delta\sigma\) be an admissible calibration signal, and assume the
packages FB-A--FB-C and DF-A(1)--DF-A(3). Then for every sufficiently
small \(\eta\in[0,\eta_*)\), the closed-loop datum
\(\mathfrak C_\eta\) of Definition~\ref{def:closed} is well-defined.
Equivalently, the diagram
\[
S_{\mathrm{source}}
\longrightarrow
 \mathcal O_B
\longrightarrow
\delta\sigma
\longrightarrow
S_{\mathrm{source}}^{[\eta]}
\longrightarrow
 \mathcal O_B^{[\eta]}
 \longrightarrow
 (\mathcal A_c^{[\eta]},\omega^{[\eta]})
\]
is conditionally closed at the one-step level.
\end{theorem}

\begin{proof}
The baseline arrow
\[
S_{\mathrm{source}}\to \mathcal O_B \to (\mathcal A_c,\omega)
\]
is the chosen Foundation~I chain. The arrow
\[
\mathcal O_B \to\delta\sigma
\]
is the imported T-DOME~III calibration interface. By
Theorem~\ref{thm:preserve}, the feedback step
\[
\delta\sigma \to S_{\mathrm{source}}^{[\eta]}
\]
is well-defined at the source-datum level. By
Proposition~\ref{prop:chain}, the same conditional Foundation~I chain
can then be run again on \(S_{\mathrm{source}}^{[\eta]}\), yielding
\[
S_{\mathrm{source}}^{[\eta]}
\to
\mathcal O_B^{[\eta]}
\to
(\mathcal A_c^{[\eta]},\omega^{[\eta]}).
\]
Thus every component of \(\mathfrak C_\eta\) is defined. The modified
deficit-floor comparison is given by Theorem~\ref{thm:deficit}, so the
closed loop is not only diagrammatically defined but also quantitatively
controlled at the Paper~III lower-floor level.
\end{proof}

\begin{corollary}[One-step closed-loop map]
\label{cor:map}
Under the hypotheses of Theorem~\ref{thm:closed}, there is a well-defined
map
\[
\mathscr L_\eta:
\bigl(S_{\mathrm{source}},\delta\sigma\bigr)
\longmapsto
\bigl(
S_{\mathrm{source}}^{[\eta]},
\mathcal O_B^{[\eta]},
(\mathcal A_c^{[\eta]},\omega^{[\eta]}),
\underline{\varepsilon}_{\mathrm{loc}}^{[\eta]}(n)
\bigr).
\]
\end{corollary}

\begin{proof}
This is just a repackaging of Theorem~\ref{thm:closed} and
Theorem~\ref{thm:deficit}.
\end{proof}

\begin{remark}[What Theorem~\ref{thm:closed} does not prove]
The theorem is a one-step well-posedness statement only. It does not
claim that repeated application of \(\mathscr L_\eta\) converges, that a
fixed point exists, that a fixed point would be unique, or that the
closed loop is thermodynamically optimal. It also does not address the
multi-source feedback problem in which several manifest sectors update a
shared ambient semigroup simultaneously.
\end{remark}

\section{Open Problems}\label{sec:open}

The present note closes only the smallest honest loop. The actual
programme still leaves several substantive gaps.

\begin{enumerate}
\item \textbf{Constructive origin of the feedback lift.} The paper takes
  the source-level semigroup lift
  \(\mathscr F^{[\delta\sigma]}\) as named input. A future result would
  need to derive that lift from the calibration signal itself, rather
  than postulate it.
\item \textbf{Generator-level perturbation theory.} The present paper
  works with an exact auxiliary semigroup because it preserves complete
  positivity and the semigroup law cleanly. A sharper generator-level
  perturbation theorem on \(W^*\)-algebras, closer to the existing
  structural theory of dynamical semigroups on \(W^*\)-algebras
  \cite{FrigerioVerri1982}, would have to control densely defined
  predual generators, domain stability, closability, and resolvent
  bounds in a regime where pointwise \(\sigma\)-weak continuity is the
  natural topology and norm continuity is usually unavailable. Naive
  first-order generator perturbations may fail to preserve complete
  positivity, unitality, or even semigroup well-posedness.
\item \textbf{Weakening the commutation input.} FB-C requires strict
  algebra-level commutation between the feedback lift and the source
  semigroup. A weaker theorem based on approximate commutation,
  asymptotic commutation, or commutation only on the relevant stable
  subalgebras would be more flexible and physically more natural, but it
  would need a replacement for the exact time-ordering argument used in
  Section~\ref{sec:preservation}.
\item \textbf{Quantitative coupling thresholds.} The small-coupling
  window \([0,\eta_*)\) is left abstract. The programme still lacks an
  explicit estimate of \(\eta_*\) in terms of Lyapunov constants,
  semigroup norms, and the local BKM/Fisher geometry.
\item \textbf{Persistence of the Foundation~I chain.} DF-A(1) is a real
  interface burden. One needs a theorem stating when persistent
  projections, foreground frames, and stable-hull inputs are robust
  under source-side feedback.
\item \textbf{Actual minimax deficit transport.} Theorem~\ref{thm:deficit}
  controls only the Paper~III lower floor. A full transport theorem for
  the actual minimax optimizer would be needed to compare
  \(\varepsilon_{\mathrm{loc}}^{[\eta]}(n)\) and
  \(\varepsilon_{\mathrm{loc}}(n)\) without a saturation hypothesis.
\item \textbf{Iterated perfuming and fixed points.} Once
  \(\mathscr L_\eta\) is applied repeatedly, one faces an entirely new
  dynamical question: do the semigroups converge, cycle, or wander? The
  present paper does not address that long-time behavior.
\item \textbf{Thermodynamic cost of the source-level loop.} T-DOME~III
  quantified the thermodynamic cost of the calibration loop at the
  manifest/ego level. Foundation~III has not yet integrated that cost
  with the source-side semigroup perturbation.
\item \textbf{Multi-source feedback.} Foundation~II handled the
  two-source shared-sector problem at inclusion level, but it did not
  address simultaneous feedback from multiple manifest sectors onto one
  common ambient source dynamics. Cross-source feedback remains open.
\item \textbf{Type~III modular obstacles.} Even if a feedback lift is
  available, one still faces the usual Type~III issues: modular
  invariance of the relevant subalgebras, possible non-uniqueness of
  conditional expectations, central-decomposition effects in
  non-factorial regimes, and the absence of a canonical bounded
  algebra-side generator picture that could be perturbed by finite-rank
  feedback data.
\end{enumerate}

\section{Conditional Input Ledger}\label{sec:ledger}

The theorem chain depends on the following named inputs.

\begin{enumerate}
\item \textbf{FI-A: Foundation~I source-to-manifest interface.}
  The single-source conditional chain from a source datum to a manifest
  pair is available in the sense restated in
  Appendix~\ref{app:interfaces}.
\item \textbf{TD3-A: admissible calibration signal.}
  A one-step natural-gradient update signal \(\delta\sigma\) exists and
  lies inside the T-DOME~III Lyapunov neighborhood.
\item \textbf{FB-A: CP source lift.}
  The calibration signal admits a pointwise \(\sigma\)-weakly continuous
  semigroup lift by normal completely positive unital maps on
  \(\mathcal A_U\).
\item \textbf{FB-B: small-signal differentiability.}
  The lifted feedback semigroup is differentiable at zero coupling so
  that the first-order map \(\Gamma\) is defined; its normalization
  \(\mathfrak G_{\delta\sigma}(\mathbf 1)=0\) is inherited from
  unitality.
\item \textbf{FB-C: chain compatibility.}
  The feedback lift commutes with the ambient source semigroup and
  preserves the named structural inputs needed to rerun the chosen
  Foundation~I chain.
\item \textbf{DF-A(1): feedback-stable Foundation~I chain.}
  The chosen Foundation~I conditional inputs remain valid for
  \(S_{\mathrm{source}}^{[\eta]}\) on a nontrivial coupling window.
\item \textbf{DF-A(2): manifest Fisher control.}
  The modified manifest-side Fisher sensitivity and ego-rigidity terms
  satisfy the comparison inequalities controlled by
  \(\alpha_F(\eta)\) and \(\beta_{\mathrm{ego}}(\eta)\).
\item \textbf{DF-A(3): local Paper~III control.}
  The modified manifest pair remains inside the local theorem domain of
  Paper~III, with quadratic-loss parameter \(\rho(\eta)<\tfrac12\).
\end{enumerate}

\appendix
\section{Imported Interfaces}\label{app:interfaces}

This appendix records the external results used by the body in a
self-contained theorem-ready form. The statements are imported
interfaces, not re-proved results.

\begin{remark}[What is and is not imported]
Foundation~III imports theorem-level interfaces only from
Foundation~I, T-DOME~III, and Paper~III. Foundation~II is cited in the
Introduction solely to locate the later multi-source feedback problem
and does not enter any proof or named input in the present paper.
\end{remark}

\begin{proposition}[Imported Foundation~I interface]\label{prop:A1}
Let
\[
S_{\mathrm{source}}
=
(\mathcal H_U,\mathcal A_U,\Omega,\{\mathcal E_t\}_{t\ge0})
\]
be a source datum, where
\((\mathcal H_U,\mathcal A_U,\Omega)\) is a Type~III\(_1\) standard-form
von Neumann algebra and
\(\{\mathcal E_t\}_{t\ge0}\) is a pointwise \(\sigma\)-weakly
continuous semigroup of normal completely positive unital maps on
\(\mathcal A_U\). Suppose:
\begin{enumerate}
\item there exists a persistent Nakajima--Zwanzig projection
  \(\mathcal P\) whose projected predual dynamics admits a memory kernel
  \(\mathcal K_*(t,s)\)~\cite{Nakajima1958,Zwanzig1960,BreuerPetruccione2002};
\item the resulting record process exceeds its processing budget and
  forces a compressed foreground frame \(\mathfrak F\);
\item the foreground transfer/stable-hull construction closes so as to
  produce a von Neumann algebra \(\mathcal A_c\subseteq\mathcal A_U\)
  with faithful restricted state
  \(\omega=\langle\Omega,\cdot\,\Omega\rangle|_{\mathcal A_c}\);
\item the admissibility input needed to feed Paper~I~\cite{PaperI} is
  available.
\end{enumerate}
Then the conditional chain
\[
S_{\mathrm{source}}
\to
\mathcal O_A
\to
\mathcal O_B
\to
(\mathcal A_c,\omega)
\]
is defined.
\end{proposition}

\begin{remark}
Foundation~III uses only the existence of the manifest pair produced by
Proposition~\ref{prop:A1}. It does not revisit the internal stable-hull
construction of Foundation~I. The source-datum continuity imported here
is pointwise \(\sigma\)-weak continuity, matching Foundation~I v2. The
published Foundation~I v1 uses the stronger phrase ``strongly
continuous''; the downstream Foundation~I chain used here does not
depend on that stronger condition.
\end{remark}

\begin{proposition}[Imported T-DOME~III calibration interface]
\label{prop:A2}
Let \(\mathcal O_B\) be a Stage~B output in the sense of
Foundation~I. In the local calibration regime of T-DOME~III there exist:
\begin{enumerate}
\item a natural-gradient update signal
  \(\delta\sigma\in T_\sigma\Sigma\);
\item a Lyapunov neighborhood \(\mathcal U_{\mathrm{cal}}\) on which the
  one-step update stays controlled;
\item a prior Fisher term \(\mathcal I_{\mathrm{ego}}\ge0\), interpreted
  as ego rigidity, entering the van Trees bound for local drift
  estimation.
\end{enumerate}
These data determine the one-step manifest-side calibration step used in
the body of the present paper.
\end{proposition}

\begin{remark}
Proposition~\ref{prop:A2} is a one-step interface only. It does not
assert global convergence of repeated calibration updates.
\end{remark}

\begin{proposition}[Imported Paper~III local deficit floor]
\label{prop:A3}
Let \((\mathcal A_c,\omega)\) be a manifest accessible pair and let
\((\mathcal A_{\mathrm{obs}},\omega_{\mathrm{obs}},\mathfrak M_n)\) be a
finite observer carried by \((\mathcal A_c,\omega)\). Assume that the
local Paper~III hypotheses hold on a BKM neighborhood around
\(\omega_{\mathrm{obs}}\), with quadratic-loss parameter
\(\rho\in[0,\tfrac12)\). Then the local self-description deficit
\(\varepsilon_{\mathrm{loc}}(n)\) satisfies
\[
\varepsilon_{\mathrm{loc}}(n)
\ge
\frac{(1-2\rho)\,\mathcal I_F(\omega_{\mathrm{obs}})}
{2\bigl(
\mathcal I_F(\omega_{\mathrm{obs}})
+
\mathcal I_{\mathrm{ego}}
\bigr)}.
\]
\end{proposition}

\begin{remark}
The present paper imports Proposition~\ref{prop:A3} as a manifest-side
comparison theorem. The parameter \(\rho\) is written with a new symbol
here to avoid clash with the source-level coupling strength
\(\eta\) used in the body.
\end{remark}

\newpage
\addcontentsline{toc}{section}{References}

\begin{thebibliography}{99}
\sloppy

\bibitem{FoundationI}
S.~Liu,
\emph{Foundation I: From Source to Manifestation: A Conditional Same-Root Program for Geometry and Modular Flow},
Consciousness Series (2026),
Zenodo, DOI:~10.5281/zenodo.19627786.

\bibitem{FoundationII}
S.~Liu,
\emph{Foundation II: Shared Manifestation: A Conditional Two-Source Program for Shared Accessible Sectors},
Consciousness Series (2026),
Zenodo, DOI:~10.5281/zenodo.19641372.

\bibitem{PaperI}
S.~Liu,
\emph{The Inseparability of Geometry and Modular Flow: A Structural Inseparability Theorem and Its Implications for the Hard Problem},
Consciousness Trilogy, Paper~I (2026),
Zenodo, DOI:~10.5281/zenodo.19607387.

\bibitem{PaperII}
S.~Liu,
\emph{Higher-Dimensional Consistency of Geometry and Modular Flow: Small-Ball Modular Data and Linearized Gravity},
Consciousness Trilogy, Paper~II (2026),
Zenodo, DOI:~10.5281/zenodo.19607543.

\bibitem{PaperIII}
S.~Liu,
\emph{The Structural Horizon of Self-Description: Fisher Bounds, Small-Ball Remainders, and the Quantitative Explanatory Gap},
Consciousness Trilogy, Paper~III (2026),
Zenodo, DOI:~10.5281/zenodo.19607724.

\bibitem{TDOMEIII}
S.~Liu,
\emph{Fisher Information Geometry and the Thermodynamic Cost of Self-Referential Calibration},
Zenodo (2026), DOI:~10.5281/zenodo.18591771.

\bibitem{BreuerPetruccione2002}
H.-P.~Breuer and F.~Petruccione,
\emph{The Theory of Open Quantum Systems},
Oxford University Press (2002).

\bibitem{Nakajima1958}
S.~Nakajima,
\emph{On quantum theory of transport phenomena},
Prog.\ Theor.\ Phys.\ \textbf{20}, 948--959 (1958),
DOI:~10.1143/PTP.20.948.

\bibitem{Zwanzig1960}
R.~Zwanzig,
\emph{Ensemble method in the theory of irreversibility},
J.\ Chem.\ Phys.\ \textbf{33}, 1338--1341 (1960),
DOI:~10.1063/1.1731409.

\bibitem{FrigerioVerri1982}
A.~Frigerio and M.~Verri,
\emph{Long-time asymptotic properties of dynamical semigroups on \(W^*\)-algebras},
Math.\ Z.\ \textbf{180}, 275--286 (1982),
DOI:~10.1007/BF01318911.

\end{thebibliography}

\end{document}
